\documentclass[10pt,a4paper]{amsart}
\usepackage[margin=19mm]{geometry}
\usepackage{amsmath,amssymb,mathtools}
\usepackage[scr=rsfs,cal=euler]{mathalfa}
\usepackage{xfrac}
\usepackage[unicode,hidelinks,bookmarks=false]{hyperref}
\newcommand{\ie}{i.e.\ }
\newcommand{\cf}{cf.\ }
\newcommand{\resp}{respectively\ }
\newcommand{\Subsection}{\subsection}
\providecommand{\nobreakdash}{\nobreak\hbox{-}}
\setlength{\emergencystretch}{2em}
\multlinegap0pt
\usepackage{amssymb}
\newtheorem{proposition}{Proposition}
\newcommand{\BV}{\mathrm{BV}}
\newcommand{\N}{\mathbb{N}}
\title{Homogenization in fractional phase transitions \\ at the critical scale}
\author{Andrea Braides, Abdilaziz Zoir Ugli Komilov}
\date{Source: arXiv:2606.08134v1 (CC BY 4.0)}
\begin{document}
\maketitle

\noindent\emph{Source and licence.} Homogenization in fractional phase transitions \\ at the critical scale. Version arXiv:2606.08134v1 (CC BY 4.0), distributed by arXiv under the Creative Commons Attribution 4.0 International licence, http://creativecommons.org/licenses/by/4.0/, as recorded in the arXiv licence metadata for this exact version and shown on the abstract page https://arxiv.org/abs/2606.08134v1. Full licence notice: \texttt{licenses/arxiv-2606.08134-CC-BY-4.0.txt}.\par\smallskip

\noindent\textit{Editorial scope.} Counted unit: the article's proposition compactness (source0.tex lines 426--428). The article's complete original proof of it is delivered with it (source0.tex lines 430--503, one \texttt{proof} environment of 74 lines). No mathematics is written, repaired or re-derived: the statement and the proof are the article's own lines, in the article's own notation, and no retained line is substituted or re-flowed.  Retained uncounted as the source-local context the proof needs: lines 407--412.  Nothing was omitted from any retained span, so the package carries no omission record.  No retained line contains a citation, so the wrapper carries no bibliography and no bibliography entry.  No retained line contains a figure environment, an image-inclusion command or drawing code, so no image is delivered and the package declares no asset dependency.  Editorial formatting only, in addition: an AMS \texttt{amsart} preamble that restates the article's own declarations and macros used by the retained lines, namely the environments \texttt{proposition} on separate counters, together with the standard packages those lines need.  The retained environments print in this document as Proposition 1 in order of appearance, whereas the article prints the same items with the numbering of its own full document.  The complete native source of the article as distributed in the arXiv e-print for this exact version is bundled unchanged as \texttt{sources/202309/source0.tex}.
\begin{proposition}\label{convergence in measure}
   Let $(u_{\varepsilon})_{\varepsilon}$ be a sequence and let $C>0$ be such that $\sup_{\varepsilon>0}F_{\varepsilon}(u_{\varepsilon})\le C<\infty$. Then, for every $\tau>0$ we have
\[
\lim_{\varepsilon\to0}|\{t\in(0,1):\,\operatorname{dist}(u_{\varepsilon}(t),\{-1,1\})\ge\tau\}|=0
\]
\end{proposition}

\begin{proposition}\label{compactness}
Let $(u_{\varepsilon})_{\varepsilon}$ be a sequence and let $C>0$ be such that $\sup_{\varepsilon>0}F_{\varepsilon}(u_{\varepsilon})\le C<\infty$. Then, there exists a subsequence $(u_{\varepsilon_{i}})$ such that $u_{\varepsilon_{i}}$ converges in $L^{2}(0,1)$ to some $u\in\BV((0,1);\{-1,1\})$.
\end{proposition}

\begin{proof}
Let $\{I_{k}(\varepsilon)\}_{k=1}^{n(\varepsilon)+1}$ be the partition defined above. With respect to this partition, we define the sequence of functions $(v^{\varepsilon,\tau})$ as follows:
\begin{equation*}
v^{\varepsilon,\tau}(t) =
\begin{cases}
-1 & \hbox{if}\,\, t \in I_{k}({\varepsilon})\,\, k \in P_{1}(\varepsilon), \\
1 & \hbox{if}\,\, t \in I_{k}({\varepsilon})\,\, k \in P_{2}(\varepsilon)\cup P_{3}(\varepsilon)\cup\{n(\varepsilon)+1\}.
\end{cases}
\end{equation*}
We now claim that, for every fixed $\tau \in (0,1/8)$, the sequence $(v^{\varepsilon,\tau})$ admits a subsequence that converges almost everywhere in $(0,1)$ to a function $v^{\tau}\in BV((0,1);\{-1,1\})$.

The uniform bound on $\#(P_{3}(\varepsilon)\cup \bar{P}_{3}(\varepsilon))$ implies that the number of jump points of $v^{\varepsilon,\tau}$ is uniformly bounded with respect to $\varepsilon$ and $\tau$. Hence, passing to a subsequence if necessary, we may assume that the number of jump points is constant. Let $(v^{\varepsilon_{j},\tau})$ be a subsequence of $(v^{\varepsilon,\tau})$ such that the number of jump points of $(v^{\varepsilon_{j},\tau})$ is constant and equal to $N({\tau})$ for some $N({\tau}) \in \mathbb{N}$. For each $j$, let $\{t_n(\varepsilon_j)\}_{n=1}^{N(\tau)}$ denote the jump points of $v^{\varepsilon_j,\tau}$. By passing to a further subsequence if necessary, we may assume that
$t_n(\varepsilon_j)\to t_n$ as $j\to\infty$ for all $n=1,\ldots,N(\tau)$, $v^{\varepsilon_j,\tau}\to v^{\tau}$ as $j\to\infty$ pointwise almost everywhere for some $v^{\tau}\in BV((0,1);\{-1,1\})$, and the jump set of $v^{\tau}$ is contained in $\{t_1,\ldots,t_{N(\tau)}\}$. 
Analogously, after passing to a subsequence if necessary, we may assume that $v^{\tau}\to v$ pointwise almost everywhere as $\tau\to 0$, for some $v\in BV((0,1);\{-1,1\})$. Hence, after passing to a further subsequence and letting first $\varepsilon\to0$ and then $\tau\to0$, we may assume that $v^{\varepsilon,\tau}$ converges pointwise almost everywhere \ to some $v \in BV((0,1);\{-1,1\})$. In particular, both $v^{\varepsilon,\tau}$ and $v^{\tau}$ are uniformly bounded in $L^{\infty}(0,1)$. Therefore, by the dominated convergence theorem, almost everywhere convergence implies convergence in $L^{2}(0,1)$. Finally, by a diagonal argument, we may extract a sequence $\varepsilon_i\to0$ and a subsequence $\tau_k\to0$ such that, for every $k\in\N$, $v^{\varepsilon_i,\tau_k}\to v^{\tau_k}$ as $i\to\infty$, while $v^{\tau_k}\to v$ as $k\to\infty$. Moreover, the sequence $\varepsilon_i$ is independent of $k$. We now show that $u_{\varepsilon_i}\to v$ as $i\to\infty$ in $L^{2}(0,1)$.

%Finally, by a diagonal argument, there exists a function $\varepsilon(\tau)\to0$ as $\tau\to0$ such that
%$v^{\tau}_{\varepsilon(\tau)} \to u$ pointwise almost everywhere  in $(0,1)$, with $u \in BV((0,1);\{-1,1\})$.
We now claim that for every $\tau \in (0,1/8)$, the following estimate holds:
\begin{equation}\label{36tau}
    \int_{0}^{1}\left|u_{\varepsilon}(t)-v^{\varepsilon,\tau}(t)\right|^2dt\le36\tau+o(1),
\end{equation}
as $\varepsilon\to0$. Indeed, we decompose 
\begin{equation*}
    \int_{0}^{1}\left|u_{\varepsilon}(t)-v^{\varepsilon,\tau}(t)\right|^2dt=\int_{\{\left|u_{\varepsilon}\right|>1+\tau\}}\left|u_{\varepsilon}(t)-v^{\varepsilon,\tau}(t)\right|^2dt+\int_{\{\left|u_{\varepsilon}\right|\le1+\tau\}}\left|u_{\varepsilon}(t)-v^{\varepsilon,\tau}(t)\right|^2dt.
\end{equation*}
We claim that 
\[
\lim_{\varepsilon\to0}\int_{\{\left|u_{\varepsilon}\right|>1+\tau\}}\left|u_{\varepsilon}(t)-v^{\varepsilon,\tau}(t)\right|^2dt=0.
\]
By Hypothesis {\rm (A2)}, we obtain the following upper bound:
\begin{align*}
    \int_{\{\left|u_{\varepsilon}\right|>1+\tau\}}\left|u_{\varepsilon}(t)-v^{\varepsilon,\tau}(t)\right|^2dt\le2\int_{\{\left|u_{\varepsilon}\right|>1+\tau\}}(u^{2}_{\varepsilon}(t)+1)dt\\
\le\frac{2}{c_{1}}\int_{0}^{1}W(u_{\varepsilon}(t))dt+2\left(\frac{c_{2}}{c_{1}}+1\right)|\{\left|u_{\varepsilon}\right|>1+\tau\}|.
\end{align*}
Consequently, by Proposition \ref{convergence in measure}
\[
\lim_{\varepsilon\to0}\int_{\{\left|u_{\varepsilon}\right|>1+\tau\}}\left|u_{\varepsilon}(t)-v^{\varepsilon,\tau}(t)\right|^2dt\le\lim_{\varepsilon\to0}\left(\frac{2C\varepsilon|\log{\varepsilon}|}{c_{1}}+2\left(\frac{c_{2}}{c_{1}}+1\right)|\{\left|u_{\varepsilon}\right|>1+\tau\}|\right)=0.
\]
Now we will show that
\[
\int_{\{\left|u_{\varepsilon}\right|\le1+\tau\}}\left|u_{\varepsilon}(t)-v^{\varepsilon,\tau}(t)\right|^2dt\le 36\tau+o(1),
\]
as $\varepsilon\to0$. We define the sets $J_1(\varepsilon)$, $J_2(\varepsilon)$ and $J_3(\varepsilon)$ as follows:
\begin{equation*}
    J_{1}(\varepsilon):=\bigcup_{k\in P_{1}(\varepsilon)}I_{k}(\varepsilon)\bigcap \{\left|u_{\varepsilon}\right|\le{1+\tau\}},\,\,J_{2}(\varepsilon):=\bigcup_{k\in P_{2}(\varepsilon)}I_{k}(\varepsilon)\bigcap \{\left|u_{\varepsilon}\right|\le{1+\tau\}},
\end{equation*}
\begin{equation*}
    J_{3}(\varepsilon):=I_{n(\varepsilon)+1}\bigcup_{k\in P_{3}(\varepsilon)}I_{k}(\varepsilon)\bigcap \{\left|u_{\varepsilon}\right|\le{1+\tau\}}.
\end{equation*}
Then, we have 
\begin{multline*}
    \int_{\{\left|u_{\varepsilon}\right|\le1+\tau\}}\left|u_{\varepsilon}(t)-v^{\varepsilon,\tau}(t)\right|^2dt= \int_{J_{1}(\varepsilon)}\left|u_{\varepsilon}(t)+1\right|^2dt+\int_{J_{2}(\varepsilon)}\left|u_{\varepsilon}(t)-1\right|^2dt\\+\int_{J_{3}(\varepsilon)}\left|u_{\varepsilon}(t)-1\right|^2dt\le\int_{J_{1}(\varepsilon)}\left|u_{\varepsilon}(t)+1\right|^2dt+\int_{J_{2}(\varepsilon)}\left|u_{\varepsilon}(t)-1\right|^2dt\\+(2+\tau)^{2}(\#P_{3}(\varepsilon)+1)\theta(\varepsilon).
\end{multline*}
Then, by the uniform bound on $\#P_3(\varepsilon)$, we obtain
\begin{equation*}
\int_{\{\left|u_{\varepsilon}\right|\le1+\tau\}}
\left|u_{\varepsilon}(t)-v^{\varepsilon,\tau}(t)\right|^2dt
\le\int_{J_{1}(\varepsilon)}
\left|u_{\varepsilon}(t)+1\right|^2dt+\int_{J_{2}(\varepsilon)}
\left|u_{\varepsilon}(t)-1\right|^2dt+o(1),
\end{equation*}
as $\varepsilon\to0$. We claim that 
\begin{align*}
     \int_{J_{1}(\varepsilon)}\left|u_{\varepsilon}(t)+1\right|^2dt\le18\tau,\,\int_{J_{2}(\varepsilon)}\left|u_{\varepsilon}(t)-1\right|^2dt\le18\tau.
\end{align*}
We give the proof for the first estimate; the second one follows similarly
 \begin{multline*}
     \int_{J_{1}(\varepsilon)}\left|u_{\varepsilon}(t)+1\right|^2dt=\int_{\cup_{k\in P_{1}}A^{c}_k(\varepsilon)\cap\{\left|u_{\varepsilon}\right|\le\tau+1\}}\left|u_{\varepsilon}(t)+1\right|^2dt\\+\int_{\cup_{k\in P_{1}(\varepsilon)}A_k(\varepsilon)\cap\{\left|u_{\varepsilon}\right|\le\tau+1\}}\left|u_{\varepsilon}(t)+1\right|^2dt\le2\tau(2+\tau)^{2}+\tau^{2}\le18\tau.
 \end{multline*}
 The above estimates yield the claimed result. Then, by \eqref{36tau} and the convergence of $v^{\varepsilon_i,\tau_k}\to v^{\tau_k}$ in $L^{2}(0,1)$ as $i\to\infty$, we obtain
\begin{multline*} \limsup_{i\to\infty}||u_{\varepsilon_{i}}-v|| \le\limsup_{i\to\infty}||u_{\varepsilon_{i}}-v^{\varepsilon_{i},\tau_{k}}||_{L^{2}(0,1)}+\limsup_{i\to\infty}||v^{\varepsilon_{i},\tau_{k}}-v^{\tau_{k}}||_{L^{2}(0,1)}\\+||v^{\tau_{k}}-v||_{L^{2}(0,1)}\le6\sqrt{\tau_{k}}+||v^{\tau_{k}}-v||_{L^{2}(0,1)} 
\end{multline*}
for every $k \in \mathbb{N}$. By arbitrariness of $k$, letting $k \to \infty$ yields $\lim_{i\to\infty} \|u_{\varepsilon_i}-v\|_{L^2(0,1)} = 0,$ that is, $u_{\varepsilon_i} \to v$ strongly in $L^2(0,1)$. Finally, passing to a further subsequence, we may assume that the convergence holds almost everywhere in $(0,1)$. This concludes the proof of the compactness result.
\end{proof}

\end{document}
